% VAMOS — Title Page (with author identifiers)
% Submitted to: Neurocomputing (Elsevier)
% Separate title-page source; main.tex is the complete manuscript with author details.
% =============================================================================
\documentclass[final,5p,times,twocolumn]{elsarticle}

\usepackage{graphicx}
\usepackage{amsmath}
\usepackage{amssymb}
\usepackage{xcolor}
\usepackage[hidelinks]{hyperref}
\hypersetup{hypertexnames=false}
% Editorial changes relative to repository commit 2aa37a0.
% Keep revision markup separate from the scientific content and BibTeX records.
\usepackage{etoolbox}
\DeclareRobustCommand{\revision}[1]{\textcolor{red}{#1}}
\pdfstringdefDisableCommands{\def\revision#1{#1}\def\corref#1{}}

% Mark the two references requested for this revision without editing a .bbl.
\newcommand{\markrevisionreferences}{%
  \let\revisionoriginalbibitem\bibitem
  \renewcommand{\bibitem}[2][]{%
    \par
    \ifstrequal{##2}{HMM+22}{\color{red}}{%
      \ifstrequal{##2}{TCZ+17}{\color{red}}{\color{black}}}%
    \ifstrempty{##1}{\revisionoriginalbibitem{##2}}{%
      \revisionoriginalbibitem[##1]{##2}}}}

% The regenerated source table has revised repository links and source labels.
\makeatletter
\newenvironment{revisiontable}{%
  \begingroup
  \let\revisionoriginalfloatboxreset\@floatboxreset
  \def\@floatboxreset{\revisionoriginalfloatboxreset\color{red}}%
}{\endgroup}
\makeatother


\newcommand{\VAMOS}{VAMOS}
\journal{\revision{Neurocomputing}}
\biboptions{sort&compress}

\begin{document}

\begin{frontmatter}

\title{VAMOS: A Framework for Fast, Configurable and Accessible Multiobjective Optimization in Python\tnoteref{mytitlenote}}

\tnotetext[mytitlenote]{This work has been partially supported by the Spanish
Ministerio de Ciencia e Innovaci\'on (MCIN/AEI/10.13039/501100011033) under
grant refs.\ PID2024-156045NB-I00, TSI-100930-2023-3 PID2021-125709OA-C22, and
PID2024-155363OB-C41 and by ERDF A way of making Europe; and Comunidad
Aut\'onoma de Madrid with grant ref.\ TEC-2024/COM-404 and by grant
DGP\_PIDI\_2024\_01174, Junta de Andaluc\'ia. J.\ Del Ser acknowledges funding
support from the Basque Government through the consolidated research group
MATHMODE (IT1866-26).}

\cortext[cor]{Corresponding author: Nicol\'{a}s R. Uribe
(nicolas.rodriguez@urjc.es)}

\author[URJC,cor]{Nicol\'{a}s R. Uribe}
\ead{nicolas.rodriguez@urjc.es}

\author[URJC]{Alberto Herr\'{a}n}
\ead{alberto.herran@urjc.es}

\author[LCC,ITIS]{Antonio J. Nebro}
\ead{ajnebro@uma.es}

\author[tec,upv]{Javier {Del Ser}}
\ead{javier.delser@tecnalia.com}

\author[URJC,cor]{J. Manuel Colmenar}
\ead{josemanuel.colmenar@urjc.es}

\address[URJC]{Dept.\ Computer Sciences, Universidad Rey Juan Carlos,
C/. Tulip\'{a}n, s/n, M\'{o}stoles, 28933 (Madrid), Spain}

\address[LCC]{Dept.\ de Lenguajes y Ciencias de la Computaci\'{o}n, ITIS
Software, University of M\'{a}laga, ETSI Inform\'{a}tica, Campus de Teatinos,
29071 (M\'{a}laga), Spain}

\address[ITIS]{ITIS Software, Ada Byron Research Building,
C/. Arquitecto Francisco Pe\~{n}alosa, 18, University of M\'{a}laga, 29071
(M\'{a}laga), Spain}

\address[tec]{TECNALIA, 48160 Derio, Spain}

\address[upv]{University of the Basque Country (UPV/EHU), 48013 Bilbao, Spain}

\begin{abstract}
Python is widely used for empirical research in multi-objective optimization due
to its accessibility and strong scientific ecosystem. However, most Python
frameworks store populations as per-solution objects and keep critical
computations in the interpreter, which prevents vectorization and significantly
limits performance. This design gap makes Python-based metaheuristics much
slower than their counterparts in compiled languages. In this paper we introduce
VAMOS (Vectorized Architecture for Multiobjective Optimization Suite), a
vectorized Python framework for multi-objective evolutionary algorithms (MOEAs).
VAMOS maintains all population data in dense numerical arrays and dispatches
hot-path operations to interchangeable numerical backends, decoupling algorithm
logic from the underlying compute substrate. The framework implements eight
MOEAs spanning dominance-based, decomposition-based, and indicator-based
paradigms, and exposes rich component-level configurability, including
interchangeable operators, external archives, and fine-grained parameter
settings. VAMOS also integrates an automated algorithm-configuration module and
provides a browser-based graphical interface to simplify use by non-experts. The
results show that VAMOS achieves \revision{levels} of solution quality \revision{statistically equivalent to} those \revision{obtained} by
MOEAs implemented in state-of-the-art libraries,
while reducing wall-clock time by large margins across all tested algorithms.
These gains confirm that vectorized population models and backend-agnostic
kernels at the core of VAMOS can deliver high performance in Python without
compromising algorithmic behavior.
\end{abstract}

\begin{keyword}
Evolutionary algorithms \sep Multiobjective optimization \sep Software framework
\sep Vectorization
\end{keyword}

\end{frontmatter}

\section*{Acknowledgments}

J.\ Del Ser acknowledges funding support from the Basque Government through the
consolidated research group MATHMODE (IT1866-26).

\end{document}
